\documentclass[11pt]{article}
\usepackage[T1]{fontenc}
\usepackage[utf8]{inputenc}
\usepackage{lmodern,microtype,anyfontsize}
\usepackage[a4paper,top=22mm,bottom=23mm,left=23mm,right=23mm]{geometry}
\usepackage{amsmath,amssymb,amsthm,mathtools,mathrsfs}
\usepackage{booktabs,longtable,array,tabularx,multirow,float}
\usepackage[table]{xcolor}
\usepackage{graphicx,tikz}
\usetikzlibrary{arrows.meta,positioning,calc,fit}
\usepackage{needspace,fancyhdr,titlesec,enumitem}
\usepackage[most]{tcolorbox}
\usepackage{hyperref,xurl}
\definecolor{Navy}{HTML}{244F7B}
\definecolor{Teal}{HTML}{167C80}
\definecolor{Amber}{HTML}{B77623}
\definecolor{Ice}{HTML}{F3F6FA}
\definecolor{Ink}{HTML}{202B36}
\hypersetup{colorlinks=true,linkcolor=Navy,citecolor=Teal,urlcolor=Navy,
  pdftitle={Two-round SRHT: hybrid constants and proof-chain comparison},
  pdfauthor={Method Problems -- consolidated research note},bookmarksnumbered=true}
\color{Ink}
\pagestyle{fancy}\fancyhf{}
\fancyhead[L]{\sffamily\small Two-round SRHT}
\fancyhead[R]{\sffamily\small Row-factor / energy comparison}
\fancyfoot[L]{\sffamily\footnotesize Appendix refinement\;\textbullet\;16 September 2026}
\fancyfoot[R]{\sffamily\small\thepage}
\renewcommand{\headrulewidth}{.4pt}
\setlength{\headheight}{15pt}
\setlength{\emergencystretch}{2em}
\setlength{\parskip}{3pt}
\setlength{\parindent}{0pt}
\linespread{1.04}
\allowdisplaybreaks[1]
\clubpenalty=10000\widowpenalty=10000
\titleformat{\section}{\Large\sffamily\bfseries\color{Navy}}{\thesection}{.65em}{}
\titleformat{\subsection}{\large\sffamily\bfseries\color{Navy}}{\thesubsection}{.6em}{}
\titleformat{\subsubsection}{\normalsize\sffamily\bfseries}{\thesubsubsection}{.6em}{}
\titlespacing*{\section}{0pt}{2.0ex plus 1ex minus .2ex}{1.0ex}
\titlespacing*{\subsection}{0pt}{1.8ex plus 1ex minus .2ex}{.7ex}
\numberwithin{equation}{section}
\newtheorem{theorem}{Theorem}[section]
\newtheorem{lemma}[theorem]{Lemma}
\newtheorem{proposition}[theorem]{Proposition}
\newtheorem{corollary}[theorem]{Corollary}
\theoremstyle{definition}\newtheorem{definition}[theorem]{Definition}
\theoremstyle{remark}\newtheorem{remark}[theorem]{Remark}
\newcommand{\R}{\mathbb R}\newcommand{\E}{\mathbb E}\newcommand{\Prob}{\mathbb P}
\newcommand{\tr}{\operatorname{tr}}\newcommand{\Ent}{\operatorname{Ent}}
\newcommand{\norm}[1]{\lVert#1\rVert}\newcommand{\ip}[2]{\langle#1,#2\rangle}
\newcommand{\one}{\mathbf1}\newcommand{\F}{\mathbb F}
\newcommand{\step}[1]{\par\addvspace{.65\baselineskip}\Needspace{4\baselineskip}\noindent\textbf{#1}\par\nobreak\smallskip}
\newcolumntype{Y}{>{\raggedright\arraybackslash}X}
\newcolumntype{P}[1]{>{\raggedright\arraybackslash}p{#1}}
\newtcolorbox{takeaway}[1][]{colback=Ice,colframe=Navy!50,boxrule=.45pt,
  arc=1pt,left=8pt,right=8pt,top=7pt,bottom=7pt,fonttitle=\sffamily\bfseries,
  coltitle=Navy,colbacktitle=Ice,enhanced,#1}
\setlist{nosep,leftmargin=1.5em}

\newcommand{\C}{\mathbb C}
\newcommand{\diag}{\operatorname{diag}}
\newcommand{\cyc}{\operatorname{cyc}}
\newcommand{\cF}{c_{\F,p}}
\begin{document}
\thispagestyle{empty}
{\sffamily\color{Navy}\small METHOD PROBLEMS\hfill PROOF-CHAIN MAP\par}
\vspace{3mm}
{\sffamily\bfseries\fontsize{24}{28}\selectfont Two-round SRHT\par}
\vspace{1.5mm}
{\sffamily\large Where the row-factor refinement improves the constant\par}
\vspace{2mm}
{\small Read downward. Teal shows the reference proof; blue shows the refinement.\par
Full-width boxes identify the common model, completion, and scope.\par}
\vspace{2mm}

% Match the supplied template: final-size vector drawing, anchor-based placement.
\newlength{\ProofFullWidth}
\setlength{\ProofFullWidth}{\dimexpr\textwidth-16pt\relax}
\newlength{\ProofColumnWidth}
\setlength{\ProofColumnWidth}{.5\textwidth}
\addtolength{\ProofColumnWidth}{-5mm}
\addtolength{\ProofColumnWidth}{-14pt}
\newlength{\ProofColumnOffset}
\setlength{\ProofColumnOffset}{.25\textwidth}
\addtolength{\ProofColumnOffset}{2.5mm}
\begin{center}
\begin{tikzpicture}[
  >={Latex[length=1.6mm,width=1.25mm]},
  every node/.style={font=\fontsize{10.2}{12.4}\selectfont},
  base/.style={rounded corners=1.5pt,line width=.5pt,
    align=center,inner xsep=6pt,inner ysep=5pt,outer sep=.5pt},
  shared/.style={base,draw=Navy!70,fill=Ice,text width=\ProofFullWidth},
  abox/.style={base,draw=Teal!85,fill=Teal!4,
    text width=\ProofColumnWidth,minimum height=18mm},
  bbox/.style={base,draw=Navy!75,fill=Navy!3,
    text width=\ProofColumnWidth,minimum height=18mm},
  ahead/.style={base,draw=Teal,fill=Teal,text=white,
    text width=\ProofColumnWidth,minimum height=10mm,
    font=\sffamily\bfseries\fontsize{10.4}{12.2}\selectfont},
  bhead/.style={ahead,draw=Navy,fill=Navy},
  scope/.style={shared,draw=Amber!80,fill=Amber!7},
  flow/.style={->,draw=Navy!80,line width=.6pt},
  aflow/.style={->,draw=Teal!90,line width=.6pt}
]
\node[shared,minimum height=17mm] (model) {\textbf{Same two-round sketch; same fixed subspace}\\[2pt]
       $W=AD_yBD_xU,\quad E_T=(n/M)W^*\Pi_TW-I_d$\\[1pt]
       $\Pi_T$ projects onto a uniform set of $M$ rows};
\node[ahead,below=7mm of model.south,xshift=-\ProofColumnOffset] (ah) {REFERENCE PROOF\\v11 / scalar retuning};
\node[bhead,below=7mm of model.south,xshift=\ProofColumnOffset] (bh) {REFINED PROOF\\Row-factor / energy};
\coordinate (split) at ($(model.south)+(0,-3mm)$);
\draw[aflow] (model.south)--(split)-|(ah.north);
\draw[flow] (model.south)--(split)-|(bh.north);
\node[abox,below=3.5mm of ah] (a1) {\textbf{A1. Use the original sign law}\\[2pt]
       Exact Boolean contractions;\\
       inherited coefficient-channel bounds};
\node[bbox,below=3.5mm of bh] (b1) {\textbf{B1. Keep the full sign multiplier}\\[2pt]
       $X_i=a_i+a_i^*,\quad X_i^*X_i=I$\\
       No Gaussian replacement};
\draw[aflow] (ah.south)--(a1.north);
\draw[flow] (bh.south)--(b1.north);
\node[fit=(a1)(b1),inner sep=0pt,outer sep=0pt] (row1) {};
\node[abox,below=4mm of row1.south,xshift=-\ProofColumnOffset] (a2) {\textbf{A2. Form fixed-grade Grams}\\[2pt]
       Diagonal allowance $D_q/n$;\\
       mix grades after taking the Gram};
\node[bbox,below=4mm of row1.south,xshift=\ProofColumnOffset] (b2) {\textbf{B2. Retain the rectangular row map}\\[2pt]
       Join the two $x$ transitions;\\
       apply Bessel before insertion};
\draw[aflow] (a1.south)--(a2.north);
\draw[flow] (b1.south)--(b2.north);
\node[fit=(a2)(b2),inner sep=0pt,outer sep=0pt] (row2) {};
\node[abox,below=4mm of row2.south,xshift=-\ProofColumnOffset] (a3) {\textbf{A3. Count nine Gram neighbours}\\[2pt]
       $L_q\le\min\{1,9D_q/n\}$\\
       Bound the neighbouring blocks separately};
\node[bbox,below=4mm of row2.south,xshift=\ProofColumnOffset] (b3) {\textbf{B3. Bound the row factor first}\\[2pt]
       $L_q\le\min\{1,\Lambda_q/n\}$\\
       The retained estimate gives $\Lambda_q\le3D_q$};
\draw[aflow] (a2.south)--(a3.north);
\draw[flow] (b2.south)--(b3.north);
\node[fit=(a3)(b3),inner sep=0pt,outer sep=0pt] (row3) {};
\node[abox,below=4mm of row3.south,xshift=-\ProofColumnOffset] (a4) {\textbf{A4. Bernoulli root allowance}\\[2pt]
       $6\sqrt{D_q/m}+9D_q/m$\\
       $m$ is the expected sample size};
\node[bbox,below=4mm of row3.south,xshift=\ProofColumnOffset] (b4) {\textbf{B4. Bernoulli root allowance}\\[2pt]
       $2\sqrt3\sqrt{D_q/m}+3D_q/m$\\
       The same normalization as A};
\draw[aflow] (a3.south)--(a4.north);
\draw[flow] (b3.south)--(b4.north);
\node[fit=(a4)(b4),inner sep=0pt,outer sep=0pt] (lastrow) {};
\node[shared,below=7mm of lastrow.south,minimum height=22mm] (join) {\textbf{Shared completion: positive selectors and exact cutoff}\\[2pt]
       Retain the first $q$ vacuum iterates, then compare to fixed-size sampling.\\[2pt]
       The sampling comparison costs at most $2$ in the moment root; apply Markov.};
\draw[aflow] (a4.south)--++(0,-3mm)-|([xshift=-24mm]join.north);
\draw[flow] (b4.south)--++(0,-3mm)-|([xshift=24mm]join.north);
\node[shared,below=4.5mm of join,minimum height=24mm] (result) {\textbf{Same failure event; a smaller sufficient row coefficient}\\[2pt]
       $M=\min\{n,\lceil C D_q/\varepsilon^2\rceil\},\quad
       q=\lceil\log_4(4d/\delta)\rceil$\\[3pt]
       $C:\ 1024\ \text{(printed)}\ \longrightarrow\ 646\ \text{(retuned)}
       \ \longrightarrow\ 216\ \text{(refined)}$};
\draw[flow] (join.south)--(result.north);
\node[scope,below=4.5mm of result,minimum height=17mm] (scope) {\textbf{Unchanged model, scope, and confidence order}\\[2pt]
       $D_q=\beta d+(8\beta+4\alpha)q^2$; the two sign families and\\
       the uniform fixed-size sampling law are retained};
\draw[flow] (result.south)--(scope.north);
\end{tikzpicture}
\end{center}
\clearpage

% Dedicated second-page comparison; all tables and equations here are unnumbered.
\section*{Proof and constants comparison}
{\sffamily\color{Teal}\large Same two-round law; a sharper row-factor estimate.\par}
\vspace{2mm}
Both routes use the same coherence allowance $D_q$, moment order $q$, and
uniform fixed-size sampling law. The table separates a stronger operator
inequality from merely retuning the final probability certificate.

\vspace{2mm}
\begingroup
\fontsize{9.5}{11.7}\selectfont
\setlength{\tabcolsep}{4pt}\renewcommand{\arraystretch}{1.16}
\begin{tabularx}{\textwidth}{@{}P{.24\textwidth}YY@{}}
\toprule
\textbf{Comparison} &
\textcolor{Teal}{\textbf{Reference proof}}\newline \textcolor{Teal}{Supplied v11} &
\textcolor{Navy}{\textbf{Refined proof}}\newline \textcolor{Navy}{Row-factor / energy (\S\ref{sec:improvement})}\\
\midrule
\rowcolor{Ice}
Retained object & Fixed-grade row Grams, assembled by their nine possible neighbours. &
The rectangular row factor; keep both $x$ transitions joined during $y$ creation.\\[3pt]
Local coefficient budget & $D_q/n$ on each input grade. & The same inherited channel estimates; no new local law.\\[3pt]
\rowcolor{Ice}
Subset-effect allowance & $L_q\le\min\{1,9D_q/n\}$. &
$L_q\le\min\{1,\Lambda_q/n\}$, with $\Lambda_q\le3D_q$.\\[3pt]
Bernoulli root allowance & $6\sqrt{D_q/m}+9D_q/m$ &
$2\sqrt3\sqrt{D_q/m}+3D_q/m$\\[3pt]
\rowcolor{Ice}
Fixed-size root allowance & $12\sqrt{D_q/M}+18D_q/M$ &
$4\sqrt3\sqrt{D_q/M}+6D_q/M$\\[3pt]
Row coefficient & $1024$ as printed; $646$ after scalar retuning alone. &
$216$, using the retained row-factor estimate.\\[3pt]
\rowcolor{Ice}
Unchanged scope & Arbitrary fixed isometry and coherent complex unitary pairs. &
Same model and $D_q$; the same $4q$-wise sign requirement, capped at $n$.\\[3pt]
\bottomrule
\end{tabularx}
\endgroup
{\footnotesize Root allowances exclude the common external factor $d^{1/(2q)}$.
Here $m=\rho n$ for Bernoulli sampling and $M$ is the fixed sample size.\par}

\subsection*{Sufficient row counts}
{\small Flat transforms; $\varepsilon=1/2$, $\delta=0.01$, $n=2^{30}$.
All columns use $q=\lceil\log_4(400d)\rceil$ and $D_q=d+12q^2$.
The full comparison appears in Section~\ref{sec:constants}.\par}
\vspace{1.5mm}
\begingroup
\fontsize{9.5}{11.7}\selectfont\setlength{\tabcolsep}{5pt}
\renewcommand{\arraystretch}{1.18}
\begin{tabularx}{\textwidth}{@{}r r >{\raggedleft\arraybackslash}X >{\raggedleft\arraybackslash}X >{\raggedleft\arraybackslash}X@{}}
\toprule
$d$ & $q$ & \textbf{Printed: $1024$} & \textbf{Retuned: $646$} & \textcolor{Navy}{\textbf{Refined: $216$}}\\
\midrule
\rowcolor{Ice}10 & 6 & 1\,810\,432 & 1\,142\,128 & \textbf{381\,888}\\
100 & 8 & 3\,555\,328 & 2\,242\,912 & \textbf{749\,952}\\
\rowcolor{Ice}1\,000 & 10 & 9\,011\,200 & 5\,684\,800 & \textbf{1\,900\,800}\\
10\,000 & 11 & 46\,907\,392 & 29\,591\,968 & \textbf{9\,894\,528}\\
\bottomrule
\end{tabularx}
\endgroup
{\footnotesize These are sufficient theorem prescriptions, not empirical thresholds.
The middle column prevents numerical retuning from being credited to the new contraction.\par}

\vspace{2mm}
\begin{takeaway}[title={Quantitative endpoint},top=4pt,bottom=4pt]
\small
\setlength{\abovedisplayskip}{4pt}\setlength{\belowdisplayskip}{4pt}
\[
 M=\min\{n,\lceil216D_q/\varepsilon^2\rceil\}
 \quad\Longrightarrow\quad \Prob\{\|E_T\|>\varepsilon\}\le\delta.
\]
The improvement $1024\to646$ is retuning; $646\to216$ uses the stronger
contraction. For flat transforms at confidence $0.99$, the rank-only
coefficient is $65016$ instead of $308224$.
\end{takeaway}
\clearpage
\begin{takeaway}[title=Universal improvement]
For the same two-round sketch, the same arbitrary fixed subspace, and the same coherence allowance
\[
 D_q=\beta d+(8\beta+4\alpha)q^2,
\]
we reduce the subset-effect loss from $9D_q/n$ to $3D_q/n$. The structural step bounds the \emph{rectangular row map before its Gram}: the two first-layer transitions remain joined while the second layer creates an occupation.

The resulting fixed-size sampling theorem permits
\[
 \boxed{M=\min\{n,\lceil216\varepsilon^{-2}D_q\rceil\}}
\]
with the same $q=\lceil\log_4(4d/\delta)\rceil$ as v11, instead of its coefficient $1024$. A control calculation shows that retuning the old proof alone gives $646$; thus the contraction improvement is separate from numerical retuning.

These are analytic refinements of the supplied appendices. The Rademacher signs and sampling law are unchanged. Gaussian Wick pairings without the original-law corrections would not justify the proof.
\end{takeaway}

\section{The precise theorem and the baseline}
\label{sec:model}
Let $A,B\in\C^{n\times n}$ be unitary and $U\in\C^{n\times d}$ satisfy $U^*U=I_d$. Set
\begin{equation}
 \alpha=n\max_{j,a}|A_{ja}|^2,\qquad
 \beta=n\max_{a,i}|B_{ai}|^2.
 \label{eq:coherence}
\end{equation}
Let $D_x,D_y$ be independent diagonal matrices of independent unbiased signs, and put
\[
 W=AD_yBD_xU,\qquad w_j^*=e_j^*W.
\]
For a uniform $M$-element subset $T\subseteq[n]$, independent of the signs, let $R_T$ restrict coordinates to $T$. The sketch and error are
\begin{equation}
 S=\sqrt{n/M}\,R_TAD_yBD_x,\qquad E_T=U^*S^*SU-I_d.
 \label{eq:sketch}
\end{equation}
For Walsh matrices, this is the two-round SRHT model; the theorem also covers arbitrary complex unitary pairs with their stated coherences. The term ``SRHT'' below refers to this two-round model, not to the classical single-round construction studied in \cite{TroppSRHT}.

\Needspace{16\baselineskip}
\begin{theorem}[Improved fixed-size sampling constant]
\label{thm:main}
Let $0<\varepsilon<1$ and $0<\delta<1/2$. Define
\begin{align}
 q&=\left\lceil\log_4(4d/\delta)\right\rceil,
 &D_q&=\beta d+(8\beta+4\alpha)q^2,\nonumber\\
 M&=\min\{n,\lceil216D_q/\varepsilon^2\rceil\}.
 \label{eq:newM}
\end{align}
Then every fixed isometry $U$ satisfies
\begin{equation}
 \Prob\{\norm{E_T}>\varepsilon\}\le\delta.
 \label{eq:prob}
\end{equation}
Each sign family may instead be $\min(n,4q)$-wise independent, provided the two families and the uniform subset remain mutually independent.
\end{theorem}

The supplied v11 theorem has exactly the same $D_q$, $q$, and sampling model, with $1024$ instead of $216$ \cite[coherent-unitary extension and coefficient-kernel appendix]{Framework}. We retain its stronger quadratic-order allowance rather than comparing against the earlier cubic-order statement. The proof below inherits the coefficient contractions from v11 and replaces the \emph{assembly of their norm bounds}. The local double-annihilation reflection estimate is not claimed as new.

\section{Where Wick must be adapted to the sign law}
\label{sec:law}
For unbiased independent signs,
\begin{equation}
 \E[\epsilon_i\epsilon_j\epsilon_k\epsilon_l]
 =\delta_{ij}\delta_{kl}+\delta_{ik}\delta_{jl}+\delta_{il}\delta_{jk}
       -2\one_{\{i=j=k=l\}}.
 \label{eq:signfour}
\end{equation}
The first three terms are Gaussian pairings; the final term is the sign fourth-cumulant correction. It is indispensable: a single sign has fourth moment one rather than three. Isserlis' theorem \cite{Isserlis} is the Gaussian pairing identity, not a formula for the sign law without this correction.

Use the exact Boolean occupation basis indexed by coordinate subsets. The annihilator $a_i$ removes $i$ when present, and $a_i^*$ inserts it when absent. Put $N_i=a_i^*a_i$ and $X_i=a_i+a_i^*$. On the \emph{full} occupation space,
\begin{equation}
 X_i^*X_i=I,\qquad
 a_i a_l^*=a_l^*a_i+\delta_{il}(I-2N_i).
 \label{eq:boolean}
\end{equation}
The factor $I-2N_i$ is the original-law occupation correction. At distinct sites the operators commute. The corresponding identities in the second sign family use $y$ instead of $x$.

The v11 appendix already uses \eqref{eq:boolean} to group a positive coefficient kernel with its reflection. The present refinement keeps an additional exact relation, $X_i^*X_i=I$, through the row-factor contraction instead of splitting both $x$ transitions and then reconstructing a Gram by a nine-neighbor estimate.

This is the valid adaptation of the rMPS hybrid \cite{Hybrid}: exact local contractions, including non-Gaussian corrections, followed by a retained energy bound. It is not a change of probability space or a Gaussian comparison theorem.

\section{The inherited coefficient estimates}
\label{sec:inputs}
Let $\mathcal X_r,\mathcal Y_s$ be the $r$- and $s$-subset occupation spaces. The inner product is conjugate linear in its first argument. Set
\begin{equation}
 c_i=e_i^*U,\qquad (g_{ji})_a=\sqrt n\,\overline{A_{ja}B_{ai}},
 \quad a_y(g)=\sum_a\overline{g_a}a_{y,a}.
 \label{eq:coeff}
\end{equation}
For complex coefficients, the full second-layer multiplier is
\[
 Y(g)=a_y(g)+a_y(\overline g)^*
       =\sum_a\overline{g_a}(a_{y,a}+a_{y,a}^*).
\]
The row multiplication map is
\begin{equation}
 R_j=\frac1{\sqrt n}\sum_i c_i\otimes X_i\otimes Y(g_{ji}).
 \label{eq:row}
\end{equation}
For a set $T$, collect these rows as $R_Tf=(R_jf)_{j\in T}$, retaining the row label as an orthogonal output register. The channel $R_T^{\mu\nu}$ changes the $(x,y)$ grades by $(\mu,\nu)\in\{+1,-1\}^2$.

\subsection{Bessel relations and insertion costs}
Unitarity and \eqref{eq:coherence} imply
\begin{align}
 \norm{\sum_i z_i g_{ji}}^2&\le\alpha\sum_i|z_i|^2
          &&\text{for each }j,\nonumber\\
 \sum_j|g_{ji}\rangle\langle g_{ji}|&\preceq\beta I
          &&\text{for each }i,\qquad \sum_i c_i^*c_i=I_d.
 \label{eq:bessel}
\end{align}
For example the first synthesis matrix is $\sqrt n\,\diag(\overline{A_{j,:}})\overline B$, whose norm is at most $\sqrt\alpha$. The second is $\sqrt n\,\diag(\overline{B_{:,i}})A^*$, whose norm is at most $\sqrt\beta$.

Extracting one coordinate from grade $s$ has norm $\sqrt s$: the map sends $|J\rangle$ to $\sum_{a\in J}e_a\otimes|J\setminus\{a\}\rangle$, whose squared norm is $s$. Its adjoint inserts a coordinate and has norm $\sqrt{s+1}$ on inputs of grade $s$. Hilbert-valued coefficient extensions of \eqref{eq:bessel} follow by expansion in an orthonormal basis of the passive space.

\begin{lemma}[Input-grade contraction bounds from v11]
\label{lem:inherited}
On fixed input grades $(r,s)$, with $t=|T|$, put $p=(r-1)_+$ and $v=(s-1)_+$. Then
\begin{align}
 n\big[(R_T^{++})^*R_T^{++}+(R_T^{-+})^*R_T^{-+}\big]
    &\preceq\alpha t(s+1)I,\label{eq:joint}\\
 n\norm{R_T^{+-}}^2&\le\beta(r+1)s,\label{eq:mixed}\\
 n\norm{R_T^{--}}^2&\le\beta(d+p+pv)+\alpha tv.
 \label{eq:double}
\end{align}
Double annihilation is zero if $r=0$ or $s=0$; the nonnegative right side of \eqref{eq:double} is still a permissible bound there.
\end{lemma}
\begin{proof}[Proof inputs and attribution]
These are the joint-creation, mixed-channel, and collective double-annihilation lemmas of \cite[coefficient-map section and coefficient-kernel appendix]{Framework}. Their exact roles are as follows.

For \eqref{eq:joint}, form the collected coefficient map $\sum_i\overline{g_{ji}}\otimes(c_i\otimes X_i)$. Equation~\eqref{eq:bessel} and $X_i^*X_i=I$ bound it by $\sqrt\alpha$, before an insertion of norm $\sqrt{s+1}$. At fixed $r$ its two $x$ outputs are orthogonal, giving the sum of Grams. We prove and use its stronger mixed-$x$-grade form in Section~\ref{sec:improvement}.

For \eqref{eq:mixed}, extract one $y$ coordinate, apply the collected map
\[
 (F_Tz)_{j,i}=(c_i\otimes\langle g_{ji}|)z,
 \qquad F_T^*F_T
 =\sum_i c_i^*c_i\otimes\sum_{j\in T}|g_{ji}\rangle\langle g_{ji}|
 \preceq\beta I,
\]
and insert one $x$ coordinate. The costs multiply to $\sqrt{\beta(r+1)s}$.

For \eqref{eq:double}, the output Gram on grades $(p,v)$ has blocks
\[
 \mathcal G_{jk}=\sum_{i,l}P_{il}\,
       a_{x,i}a_{x,l}^*\otimes a_y(g_{ji})a_y(g_{kl})^*,
 \qquad P=UU^*.
\]
Its exact Boolean normal ordering gives five terms $\mathcal E+\mathcal X-\mathcal F+\mathcal Y+\mathcal Z$. The four \emph{grouped} bounds proved in the appendix are
\[
 \norm{\mathcal E}\le\beta d,\quad
 \norm{\mathcal X-\mathcal F}\le\beta p,\quad
 \norm{\mathcal Y}\le\alpha tv,\quad
 \norm{\mathcal Z}\le\beta pv.
\]
The improved second term retains $\mathcal X-\mathcal F$ as a direct sum of positive coefficient lifts times $I-2N_{y,a}$; it is not obtained by dropping a negative term. The last term is an extracted two-species coefficient kernel of norm at most $\beta$, lifted at cost $pv$. These precisely specified kernel estimates are inherited, not reasserted as a new Wick lemma. Their sum proves \eqref{eq:double}.
\end{proof}

The fixed-input diagonal Gram budget obtained by adding these bounds is
\begin{equation}
 n\norm{\Pi_{r,s}R_T^*R_T\Pi_{r,s}}
 \le\beta[d+p+pv+(r+1)s]+\alpha t(v+s+1).
 \label{eq:diagonal}
\end{equation}
For $r,s\le2q$ and $t\le q$, it is at most $D_q$. The old proof extends \eqref{eq:diagonal} to all mixed input grades using positivity and nine Gram neighbors. This gives $9D_q/n$. We next replace that extension.

\section{Retain the row factor: the loss is three, not nine}
\label{sec:improvement}
Let $P_{\rm sig}$ retain both input sign grades at most $\min(2q,n)$. \emph{Do not truncate the output of $R_T$} when forming its squared norm. Then
\[
 P_{\rm sig}M_{\sum_{j\in T}w_jw_j^*}P_{\rm sig}
       =(R_TP_{\rm sig})^*(R_TP_{\rm sig}).
\]
This equality is what permits use of the full-space identity $X_i^*X_i=I$ in a finite-cutoff proof.

Define
\begin{align}
 a_q&=\alpha q(2q+1),\nonumber\\
 b_q&=\beta(d+8q^2)+\alpha q(2q-1),\nonumber\\
 \Lambda_q&=(\sqrt{a_q}+\sqrt{2b_q})^2.
 \label{eq:Lambda}
\end{align}
Notice $a_q+b_q=D_q$ and $\Lambda_q\le3D_q$ by Cauchy--Schwarz.

\begin{theorem}[Hybrid subset-effect estimate]
\label{thm:subset}
For every $q\ge1$ and $T\subseteq[n]$ with $|T|\le q$,
\begin{equation}
 \boxed{\norm{R_TP_{\rm sig}}^2\le\Lambda_q/n\le3D_q/n.}
 \label{eq:rownew}
\end{equation}
In particular the positive row effects
\[
 A_j^{\rm eff}=P_{\rm sig}M_{w_jw_j^*}P_{\rm sig}
\]
sum to the identity on the retained sign space, and their subset allowance satisfies
\begin{equation}
 \boxed{L_q:=\max_{|T|\le\min(q,n)}\norm{\sum_{j\in T}A_j^{\rm eff}}
    \le\min\{1,\Lambda_q/n\}\le\min\{1,3D_q/n\}.}
 \label{eq:Lnew}
\end{equation}
\end{theorem}
\begin{proof}
Split the row map by the direction of its $y$ transition:
\[
 R_T=\mathcal C_T+\mathcal A_T,
 \qquad \mathcal C_T=R_T^{++}+R_T^{-+},
 \qquad \mathcal A_T=R_T^{+-}+R_T^{--}.
\]
Here each sum denotes the actual operator across all input grades, not a sum of separate worst-case norms.

\step{Creation retains both $x$ transitions.}
Fix a $y$ grade $s$ but allow arbitrary mixed $x$ grades. Before the final $y$ insertion, the row-$j$ coefficient map is
\[
 G_jf=\sum_i\overline{g_{ji}}\otimes
       (c_i\otimes X_i\otimes I)f.
\]
Bessel synthesis, valid with passive Hilbert factors, gives
\[
 \norm{G_jf}^2\le\alpha\sum_i\norm{(c_i\otimes X_i\otimes I)f}^2
       =\alpha\norm f^2.
\]
The equality uses the full $x$ multiplier $X_i$, not either of its partial transitions. Insert the new $y$ coordinate at cost $\sqrt{s+1}$ and collect $t$ rows. Hence on this entire mixed-$x$ space,
\[
 n\norm{\mathcal C_T f}^2\le\alpha t(s+1)\norm f^2.
\]
Distinct input $y$ grades have distinct output $y$ grades under creation, so their ranges are orthogonal. Consequently
\begin{equation}
 \norm{\mathcal C_TP_{\rm sig}}^2
       \le\alpha q(2q+1)/n=a_q/n.
 \label{eq:creationnew}
\end{equation}

\step{Annihilation has only two incoming grades at each output.}
Decompose $f=\sum_{r,s}f_{r,s}$ by its input grades. A fixed output pair $(r',s')$ under $\mathcal A_T$ can receive contributions only from $(r'-1,s'+1)$ and $(r'+1,s'+1)$. Cauchy--Schwarz on these two vectors, followed by summing the orthogonal output grades, yields
\[
 \norm{\mathcal A_Tf}^2\le2\sum_{r,s}\norm{\mathcal A_Tf_{r,s}}^2.
\]
For a fixed input pair, the two $x$ transitions have orthogonal outputs. Lemma~\ref{lem:inherited} therefore gives
\[
 n\norm{\mathcal A_Tf_{r,s}}^2
 \le\big\{\beta[d+p+pv+(r+1)s]+\alpha tv\big\}\norm{f_{r,s}}^2.
\]
Over $r,s\le2q$ and $t\le q$, the braces are bounded by $b_q$: at the endpoint,
\[
 p+pv+(r+1)s=8q^2,\qquad tv\le q(2q-1).
\]
Thus
\begin{equation}
 \norm{\mathcal A_TP_{\rm sig}}^2\le2b_q/n.
 \label{eq:annihnew}
\end{equation}

\step{Combine only after the two retained estimates.}
The triangle inequality applied to \eqref{eq:creationnew} and \eqref{eq:annihnew} proves \eqref{eq:rownew}. The effects are positive, and their sum is the identity because $W^*W=I_d$ pointwise. Their subset sums are therefore also bounded by $I$, proving \eqref{eq:Lnew}.
\end{proof}

\subsection{Exactly where the improvement enters}
The old argument bounded the Gram by nine neighbors in the $(x,y)$ grade lattice. Bounding the row factor first already gives a generic improvement to $4D_q/n$, because each row-map output receives at most four input grade pairs. Theorem~\ref{thm:subset} improves $4$ to $3$ by \emph{not splitting the $x$ grades at all in the $y$-creation part}. The identity $X_i^*X_i=I$ then retains the relevant cancellation before any cross-grade Cauchy--Schwarz step.

The diagonal allowance $D_q$ itself is unchanged. Nor is every map with a nine-neighbor Gram guaranteed the factor three; the stronger result uses the exact two-layer factorization and the retained full $x$ multiplier.

\section{Complete the moment proof on the original selector law}
\label{sec:moments}
We include the remaining operator argument to make the constants independently checkable. The ingredients are the positive-effects and finite-sampling lemmas already in v11 \cite{Framework}.

\subsection{Positive effects and selector creation}
For arbitrary positive operators $A_j$ with $\sum_jA_j=I$, define
\[
 L_q=\max_{|T|\le q}\norm{\sum_{j\in T}A_j}.
\]
On the selector subset space of grades at most $q$, set
\[
 (Kf)_T=\sum_{j\in T}A_jf_{T\setminus\{j\}},\qquad
 (Df)_J=\Big(\sum_{j\in J}A_j\Big)f_J.
\]
Then
\begin{equation}
 \norm K\le\sqrt{L_q},\qquad \norm D\le L_q.
 \label{eq:positive}
\end{equation}
Indeed, for each output set $T$, factor the positive row through $(A_j^{1/2})_{j\in T}$ to obtain
\[
 \norm{\sum_{j\in T}A_jf_j}^2
 \le\norm{\sum_{j\in T}A_j}\sum_{j\in T}\ip{f_j}{A_jf_j}.
\]
Sum over sets $T$ of one fixed cardinality. Each input $J$ is charged $\sum_{j\notin J}A_j\preceq I$, giving the creation bound. Different input grades have orthogonal creation ranges. The occupancy bound is block diagonal.

\subsection{Bernoulli moment bound}
For independent selectors $\eta_j\sim\mathrm{Bernoulli}(\rho)$, independent of the signs, put
\[
 E_\rho=\rho^{-1}W^*\diag(\eta)W-I_d,\qquad m=\rho n.
\]
In the normalized basis $1,(\eta-\rho)/\sqrt{\rho(1-\rho)}$, multiplication by $\eta/\rho-1$ is
\[
 \begin{pmatrix}0&\sigma\\\sigma&\tau\end{pmatrix},\qquad
 \sigma=\sqrt{(1-\rho)/\rho},\quad \tau=(1-2\rho)/\rho.
\]
Apply \eqref{eq:positive} to the effects of Theorem~\ref{thm:subset}. Compress sign grades to $2q$ and selector grade to $q$. The multiplier is
\[
 M_{\rm cut}=\sigma(K+K^*)+\tau D,
\]
so
\begin{equation}
 \norm{M_{\rm cut}}\le2\sigma\sqrt{L_q}+|\tau|L_q
      \le2\sqrt{\Lambda_q/m}+\Lambda_q/m.
 \label{eq:cutnorm}
\end{equation}
One multiplication raises each sign grade by at most two and selector grade by at most one. Thus the first $q$ iterates from the vacuum agree exactly with original-law multiplication, and
\begin{align}
 \E\tr E_\rho^{2q}
 &=\sum_{a=1}^d\norm{M_{\rm cut}^q(e_a\otimes\Omega)}^2\nonumber\\
 &\le d\big[2\sqrt{\Lambda_q/m}+\Lambda_q/m\big]^{2q}\nonumber\\
 &\le d\big[2\sqrt3\sqrt{D_q/m}+3D_q/m\big]^{2q}.
 \label{eq:Bernoulli}
\end{align}
The dimension factor is $d$, not the dimension of the cutoff space. Keeping the finite sampling density in \eqref{eq:cutnorm} also permits the sharper root
\[
 2\sqrt{(1-\rho)\Lambda_q/m}+|1-2\rho|\Lambda_q/m,
\]
but no use of that additional refinement is made in the constant $216$.

\subsection{Fixed-size sampling and its exact barycenter}
\label{subsec:sampling}
Let $1\le M<n$, $\rho=M/n$, and let $H_1,\ldots,H_n$ be deterministic Hermitian matrices. For a Bernoulli subset $\mathcal B$ and a uniform $M$-subset $\mathcal T$, define
\[
 Z_{\mathcal B}=\sum_j(\one_{j\in\mathcal B}-\rho)H_j,
 \qquad Z_{\mathcal T}=\sum_j(\one_{j\in\mathcal T}-\rho)H_j.
\]
There is a coupling with
\begin{equation}
 \E[Z_{\mathcal B}\mid\mathcal T]=cZ_{\mathcal T},\qquad
 c=1-\Prob\{\mathrm{Bin}(n-1,\rho)=M\}\ge\tfrac12.
 \label{eq:barycenter}
\end{equation}
To construct it, draw $\mathcal B$ with size $K$ and remove uniform elements if $K>M$, or add uniform missing elements if $K<M$. Symmetry makes $\mathcal T$ uniform independently of $K$. With
\[
 a=\E(K-M)_+=\E(M-K)_+,\qquad v=n\rho(1-\rho),
\]
the conditional inclusion probabilities inside and outside $\mathcal T$ are $1-a/M$ and $a/(n-M)$. Their difference is $c=1-a/v$.

Let $Y\sim\mathrm{Bin}(n-1,\rho)$. The binomial identity $k\binom nk=n\binom{n-1}{k-1}$ gives
\[
 a=M[\Prob(Y\ge M)-\Prob(K\ge M+1)]
   =M(1-\rho)\Prob(Y=M)=v\Prob(Y=M).
\]
The neighboring atoms $\Prob(Y=M)$ and $\Prob(Y=M-1)$ are equal, because their ratio is $\frac{n-M}{M}\frac{\rho}{1-\rho}=1$. Therefore each is at most $1/2$, proving \eqref{eq:barycenter}.

The function $Z\mapsto\tr Z^{2q}=\norm Z_{S_{2q}}^{2q}$ is convex on Hermitian matrices. Conditional Jensen yields
\begin{equation}
 \E\tr Z_{\mathcal T}^{2q}\le c^{-2q}\E\tr Z_{\mathcal B}^{2q}
      \le2^{2q}\E\tr Z_{\mathcal B}^{2q}.
 \label{eq:sampling}
\end{equation}
This is the centered Bernoulli-to-fixed-size coupling of v11. The classical sampling literature, including \cite{GrossNesme}, concerns related comparisons; we do not substitute a without-replacement versus with-replacement statement for \eqref{eq:sampling}.

Apply \eqref{eq:sampling} conditionally on the signs to $H_j=w_jw_j^*$ and rescale by $\rho^{-1}$. With \eqref{eq:Bernoulli} this proves
\begin{equation}
 \boxed{\E\tr E_T^{2q}
 \le d\big[4\sqrt{\Lambda_q/M}+2\Lambda_q/M\big]^{2q}
 \le d\big[4\sqrt3\sqrt{D_q/M}+6D_q/M\big]^{2q}.}
 \label{eq:fixedroot}
\end{equation}
For $M=n$ the error is zero by unitarity; no transfer is needed.

\subsection{Proof of the sampling constant and moment transfer}
If $M\ge216D_q/\varepsilon^2$, the last root in \eqref{eq:fixedroot} is at most
\begin{equation}
 \varepsilon\left(\frac{\sqrt2}{3}+\frac{\varepsilon}{36}\right)
 \le\varepsilon\frac{12\sqrt2+1}{36}<\frac\varepsilon2.
 \label{eq:216check}
\end{equation}
The last strict inequality is the exact integer check $288<289$. Markov now gives $\Prob\{\norm E_T>\varepsilon\}<d4^{-q}\le\delta/4<\delta$, proving Theorem~\ref{thm:main}.

Each entry of $E_T$ has degree at most two in each sign family. A word in $\tr E_T^{2q}$ therefore involves at most $4q$ distinct coordinates of either sign family. Matching those moments transfers the full-independent proof to the stated limited-independent law. Similarly, the Bernoulli result allows $\min(n,2q)$-wise independent selectors with correct marginals, independent of both sign families. The fixed-size sample retains its uniform law; it is not assumed to have independent coordinates.

\clearpage
\section{Constants: structural improvement versus retuning}
\label{sec:constants}
The two changes should not be conflated. The row-factor estimate improves the \emph{proof inequality}; selecting $216$ then makes that inequality imply the desired failure bound. Some improvement to the printed $1024$ is possible even without the new estimate.

\begin{table}[htbp]\centering\small
\renewcommand{\arraystretch}{1.25}\setlength{\tabcolsep}{4pt}
\begin{tabularx}{\textwidth}{@{}P{.26\textwidth}Y Y@{}}
\toprule
Quantity & Supplied v11 & Retained row-factor proof\\
\midrule
Fixed-input diagonal budget & $D_q/n$ & unchanged\\
Subset-effect allowance & $9D_q/n$ & $\Lambda_q/n\le3D_q/n$\\
Bernoulli root & $6\sqrt{D_q/m}+9D_q/m$ & $2\sqrt3\sqrt{D_q/m}+3D_q/m$\\
Fixed-size root & $12\sqrt{D_q/M}+18D_q/M$ & $4\sqrt3\sqrt{D_q/M}+6D_q/M$\\
Sampling prescription coefficient & $1024$ printed; $646$ after retuning & $216$\\
Sign independence requirement & $\min(n,4q)$ per family & unchanged\\
Flat-transform allowance & $d+12q^2$ & unchanged\\
\bottomrule
\end{tabularx}
\caption{The same model, isometry, moment order, and coherence parameters. A subset allowance is also capped by $1$. The table omits this cap only to display the constants.}
\label{tab:structure}
\end{table}

\step{Retuning control.}
With the old root, a coefficient $646$ already suffices:
\[
 \frac{12}{\sqrt{646}}+\frac{18}{646}<\frac12,
 \qquad 576\cdot646=372096<372100=610^2.
\]
Thus $1024\to646$ is a scalar retuning of the existing theorem, not a hybrid contraction gain. The comparison $646\to216$ isolates the benefit at this conservative half-error target, approximately a $66.6\%$ reduction. Relative to the printed $1024$, the reduction is $78.9\%$. Neither $646$ nor $216$ is asserted optimal for all possible proof refinements or all choices of $q$.

\subsection{Concrete row prescriptions}
For flat $A,B$, let $\varepsilon=1/2$, $\delta=0.01$, and take $n=2^{30}$ so none of the displayed prescriptions is capped at $n$. All columns use the same $q=\lceil\log_4(400d)\rceil$.
\begin{table}[htbp]\centering\small
\renewcommand{\arraystretch}{1.25}\setlength{\tabcolsep}{4pt}
\begin{tabularx}{\textwidth}{@{}r r r >{\raggedleft\arraybackslash}X >{\raggedleft\arraybackslash}X >{\raggedleft\arraybackslash}X@{}}
\toprule
$d$&$q$&$D_q$&v11: $1024$&Old retuned: $646$&New: $216$\\
\midrule
10&6&442&1\,810\,432&1\,142\,128&\textbf{381\,888}\\
100&8&868&3\,555\,328&2\,242\,912&\textbf{749\,952}\\
1\,000&10&2\,200&9\,011\,200&5\,684\,800&\textbf{1\,900\,800}\\
10\,000&11&11\,452&46\,907\,392&29\,591\,968&\textbf{9\,894\,528}\\
\bottomrule
\end{tabularx}
\caption{Sufficient theorem row counts, not empirical sample requirements. The middle column prevents scalar retuning from being attributed to the new contraction.}
\label{tab:rows}
\end{table}

\Needspace{13\baselineskip}
\subsection{The fixed-confidence rank-only prescription}
The v11 estimate $D_q=d+12q^2\le301d$ at $q=\lceil\log_4(400d)\rceil$ gives the following immediate corollary:
\begin{corollary}[Flat two-round transforms, confidence $0.99$]
For $0<\varepsilon<1$, flat unitary $A,B$, and each fixed rank-$d$ isometry, taking
\begin{equation}
 M=\min\{n,\lceil65016d/\varepsilon^2\rceil\}
 \label{eq:rankonly}
\end{equation}
is sufficient for failure probability at most $0.01$. The original printed coefficient is $308224$.
\end{corollary}
For clarity, the elementary rank estimate can be checked without a numerical scan. For $q=5$, $q^2/d\le25$. If $q=k\ge6$, then $d>4^{k-1}/400$, so $q^2/d<400k^2/4^{k-1}\le400\cdot36/4^5<25$; the latter sequence decreases for $k\ge6$. Thus $1+12q^2/d\le301$, and $216\cdot301=65016$.

\section{Verification, attribution, and limits}
\label{sec:limits}
The script \path{hybrid_appendices_checks.py} verifies the integer comparisons in \eqref{eq:216check} and the $646$ retuning, recomputes all entries of Table~\ref{tab:rows}, and checks small real and complex row operators built directly from Boolean masks. The latter tests compare creation, annihilation, and full row norms against \eqref{eq:creationnew}, \eqref{eq:annihnew}, and \eqref{eq:rownew}. They are supplementary numerical checks, not substitutes for the analytic inequalities. No Monte Carlo evidence is used.

The new theorem does not reduce the quadratic confidence exponent, improve the coefficient $12$ inside $d+12q^2$, change the number of transform rounds, or alter the uniform fixed-size sampling law. The original confidence obstruction and the all-row endpoint remain relevant. The more precise $\Lambda_q$ may improve individual certificates further, but Table~\ref{tab:rows} consistently uses the simple universal envelope $3D_q$.

The original kernel and selector lemmas are attributed to the supplied version, including its already improved reflection grouping. The added step is the whole-row bound before Gram relaxation. This is a genuine strengthening of that proof's inequality, but not evidence that a Wick-only, Boolean-chaos, or other operator argument could not derive it independently. The proof is analytic, not an external review or a Lean certificate; the formalization claims associated with earlier statements must not be transferred to this refinement.

\clearpage
\begin{thebibliography}{9}\raggedright
\bibitem{Framework} Diar Heidary. \emph{Graded Multiplication Operators for Random-Matrix Moments}. Version 11, exposition and attribution revision, 16 September 2026. Supplied source \path{v11_attribution_revised.tex}. Used here: the coherent-unitary theorem; collective row-channel bounds; exact Boolean double-annihilation kernels; positive-effects lemma; centered finite-sampling coupling; finite vacuum compression. The new proof retains these inputs and improves their cross-grade assembly.
\bibitem{Isserlis} Leon Isserlis. On a formula for the product-moment coefficient of any order of a normal frequency distribution in any number of variables. \emph{Biometrika} 12(1--2), 134--139, 1918. \url{https://doi.org/10.1093/biomet/12.1-2.134}. Gaussian pairing identity; the Boolean corrections here are established from the original sign law.
\bibitem{TroppSRHT} Joel A. Tropp. Improved analysis of the subsampled randomized Hadamard transform. \emph{Advances in Adaptive Data Analysis} 3(1--2), 115--126, 2011. \url{https://arxiv.org/abs/1011.1595}. Classical single-round SRHT background, not the two-round constant theorem used as the baseline here.
\bibitem{GrossNesme} David Gross and Vincent Nesme. \emph{Note on sampling without replacing from a finite collection of matrices}. arXiv:1001.2738, 2010. \url{https://arxiv.org/abs/1001.2738}. Related finite-sampling comparisons; the specific centered Bernoulli barycenter in Section~\ref{subsec:sampling} is the stated lemma from v11 and is proved here.
\bibitem{Hybrid} Method Problems working note. \emph{Gaussian rMPS Moments: Two Proof Chains and a Wick--Energy Hybrid}. 16 September 2026. Supplied source \path{rmps_three_chains_hybrid.tex}. Motivation for exact contraction data followed by a continuation-valid energy estimate; no Gaussian-core theorem is applied to signs without its required corrections.
\end{thebibliography}
\end{document}
